\documentclass[12pt,a4paper]{article}
\usepackage[utf8]{inputenc}
\usepackage[T2A]{fontenc}
\usepackage[russian]{babel}
\usepackage{amsmath,amssymb}
\usepackage{geometry}
\usepackage{hyperref}
\geometry{margin=2.5cm}

\title{От калибровочной теории дислокаций Лазара к механизму\\
электрического кулона в модели прецессирующих магнитных\\
диполей Дроздова: структурный вывод и оценка}
\author{Совместные расчёты (Claude Sonnet 5)}
\date{\today}

\begin{document}
\maketitle
\tableofcontents

\begin{abstract}
Работа последовательно излагает три связанные теоретические конструкции:
(I) динамическую калибровочную теорию дислокаций типа T(3), предложенную
М.~Лазаром; (II) её физическую интерпретацию и применение к вычислению
сил из принципа наименьшего действия, развитую А.~Я.~Брагинским;
(III) альтернативную единую модель А.~Дроздова, в которой электрон
представляется парой прецессирующих монополей Дирака, а гравитация,
кулон и обменное взаимодействие возникают из геометрии скольжения
трёхмерной ``мыльной плёнки'' -- границы четырёхмерного интеграла
действия -- вдоль искривлённых мировых трубок частиц. На основе точной
нелинейной структуры уравнений Лазара выводится явный вид членов
самозацепления, ответственных за возможное электрическое
взаимодействие между двумя нейтральными по суммарному магнитному
заряду прецессирующими диполями, и даётся честная оценка его величины.
\end{abstract}

\section{Калибровочная теория дислокаций Лазара (T(3)-теория)}

\subsection{Общая идея}

Континуальная теория дислокаций описывает кристалл со множеством
линейных дефектов (дислокаций) как упругую среду, обладающую
дополнительными степенями свободы, связанными с несовместностью
(некомпатибельностью) деформации. Лазар строит эту теорию как
калибровочную теорию группы трёхмерных трансляций $T(3)$, по прямой
аналогии с калибровочной теорией электромагнитного поля Максвелла.

Основные состояния системы -- физическая скорость $v^a$ и упругая
дисторсия (искажение) $\vartheta^a$ материального континуума. Локальная
трансляционная симметрия требует введения двух калибровочных
потенциалов: векторной 1-формы $\varphi^a$ и векторной 0-формы $\phi^a$,
которые компенсируют нарушение инвариантности при локальных сдвигах.
Гейдж-инвариантные (наблюдаемые) величины строятся как
\begin{equation}
\vartheta^a = d\xi^a+\varphi^a,\qquad v^a=\dot\xi^a+\phi^a,
\end{equation}
где $\xi^a$ -- поле деформации (аналог поля Голдстоуна).

\subsection{Полевые напряжённости: плотность и ток дислокаций}

Из калибровочных потенциалов строятся две базовые напряжённости --
прямые аналоги магнитного и электрического полей Максвелла:
\begin{align}
T^a &= d\varphi^a \quad(\text{плотность дислокаций, аналог } \vec B),\\
I^a &= d\phi^a-\dot\varphi^a \quad(\text{дислокационный ток, аналог } \vec E).
\end{align}
Они удовлетворяют тождествам Бьянки (аналоги закона Гаусса для $\vec B$
и закона Фарадея):
\begin{equation}
dT^a=0,\qquad dI^a+\dot T^a=0.
\end{equation}
Интегрирование $T^a$ по замкнутой поверхности даёт вектор Бюргерса
дислокации -- топологический аналог квантованного магнитного заряда.

\subsection{Возбуждения и определяющие соотношения}

Для замыкания теории вводятся отклики (excitations), канонически
сопряжённые напряжённостям: псевдомоментное напряжение $H_a$ и поток
дислокационного импульса $D_a$ -- прямые аналоги $\vec H$ и $\vec D$:
\begin{equation}
H_a := -\frac{\partial L_{disl}}{\partial T^a},\qquad
D_a := \frac{\partial L_{disl}}{\partial I^a}.
\end{equation}
Дополнительно вводятся дислокационное напряжение $E_a$ и
дислокационный импульс $\pi_a$:
\begin{align}
E_a &:= e_a\lrcorner L_{disl} + (e_a\lrcorner T^b)\wedge H_b
- (e_a\lrcorner I^b) D_b,\\
\pi_a &:= (e_a\lrcorner T^b)\wedge D_b = -(e_a\lrcorner D_b)\wedge T^b.
\end{align}

\subsection{Полевые уравнения типа Янга--Миллса}

Вариация полного лагранжиана по гейдж-инвариантным переменным даёт
замкнутую систему из двенадцати уравнений:
\begin{align}
dD_a - \pi_a &= p_a, \\
dH_a - \dot D_a - E_a &= \Sigma_a,
\end{align}
где $p_a$ -- упругий импульс, а $\Sigma_a$ -- упругое напряжение
(внешние, ``элaстические'' источники). Поскольку $\pi_a$ и $E_a$ сами
выражены через $T^a, I^a, H_a, D_a$, калибровочные поля <<производят
собственные источники>>, порождая \textbf{самозацепления}. Именно
поэтому уравнения нелинейны, в отличие от линейной электродинамики
Максвелла, где поля не входят в правую часть своих же уравнений.
Из этой системы выводится динамическая сила Пича--Кёлера -- прямой
аналог силы Лоренца:
\begin{equation}
f_a = -\dot\pi_a - dE_a = \dot p_a + d\Sigma_a
= (e_a\lrcorner I^b)\,p_b + (e_a\lrcorner T^b)\wedge \Sigma_b.
\end{equation}

\subsection{Линейная (квадратичная) версия теории}

Для локального, линейного, изотропного континуума определяющие
соотношения принимают вид
\begin{equation}
H_a = a_1\,gab\,{}^\star T^b,\qquad D_a = f_1\,gab\,{}^\star I^b,
\end{equation}
с константами $a_1$ (размерность силы) и $f_1$ (размерность
масса/длина). Лагранжиан дислокаций в этом приближении билинеен:
\begin{equation}
L_{disl}=\tfrac12 I^a\wedge D_a - \tfrac12 T^a\wedge H_a.
\end{equation}
Именно эта линеаризованная версия использовалась Брагинским для
получения статической силы Пича--Кёлера между дислокациями.

\section{Применение Брагинского: сила из минимума действия}

Брагинский показал, что силы, полученные Кадичем и Эделеном в
континуальной теории дислокаций по аналогии с электродинамикой, можно
вывести непосредственно из принципа наименьшего действия для тензорного
поля взаимодействия. Ключевой результат -- центрально-симметричная
компонента силы тензорного поля взаимодействия оказывается аналогом
кулоновской силы, а вихревая компонента -- аналогом силы Лоренца
(силой Пича--Кёлера в её статической форме). Важная методологическая
деталь, отмеченная самим Брагинским: в отличие от теории Кадича--Эделена,
где локально калибруется подгруппа трансляций, в развитии Брагинского
подгруппа трансляций глобальна, а представление волновой функции --
локально; это меняет трансформационные свойства компенсирующего поля
при инверсии времени.

\section{Модель Дроздова: электрон как прецессирующий магнитный диполь}

\subsection{Онтология вакуумной плёнки}

В модели Дроздова физический вакуум -- не пустота, а упругая среда
планковской жёсткости. Граница четырёхмерного интеграла действия
образует трёхмерную гиперповерхность (<<мыльную плёнку>>) в
пространстве Минковского. Единственные фундаментальные заряды --
магнитные монополи Дирака, интерпретируемые как топологические дефекты
(дислокации) решётки вакуума -- отсюда прямая связь с калибровочной
теорией дислокаций Лазара и её интерпретацией Брагинским.

\subsection{Структура электрона}

Электрон представляется парой связанных монополей $+q_m,-q_m$,
прецессирующих вокруг общей оси с частотой $\omega$ на радиусе $a$.
Электрический заряд не является фундаментальным -- он возникает как
кинематический псевдозаряд из наклона/искривления границы
интегрирования (условие Ферми) вблизи движущихся монополей.

\subsection{Условие Ферми и продавливание плёнки}

Локальное условие ортогональности плёнки к мировой трубке частицы
(Э.~Ферми, 1923) даёт для точки $O$ со скоростью, близкой к нулю:
\begin{equation}
dt = dt_0\Big(1+\frac{\vec a\cdot\vec r}{c^2}\Big).
\end{equation}
Для протяжённого (не точечного) заряда точное решение условия
ортогональности внутри толщины трубки -- не прямая, а экспонента
$t(x)=t_0 e^{-a(x-x_0)}$, что и разрешает классический парадокс
электромагнитной массы $4/3$.

\section{Три проверенных механизма электрического кулона}

\subsection{Вариант 1 -- линейная магнитостатика (Пич--Кёлер)}

Прямое (знаковое) взаимодействие всех четырёх пар монополей двух
диполей даёт точное сокращение монопольного и дипольного членов
(теорема Гаусса: суммарный заряд каждого диполя равен нулю). Ведущий
неисчезающий вклад после усреднения по прецессии:
\begin{equation}
\langle U_{signed}\rangle\cdot R = -2q_m^2\varepsilon^2\cos\Delta\phi
+O(\varepsilon^4),\qquad \varepsilon=a/R.
\end{equation}
Настоящего кулона $R^{-1}$ этот линейный механизм не даёт --
ведущий порядок $R^{-3}$.

\subsection{Вариант 2 -- линейная плёнка Ферми}

Изотропный источник продавливания $\sigma_i=m_i(\vec W_i\cdot\vec
W_i)/c^2$ даёт истинный монополь $R^{-1}$, но он не зависит от знака
заряда и от разности фаз прецессии -- то есть гравитационно-, а не
электрически-подобный эффект.

\subsection{Вариант 3 -- нелинейная самосвязь (данная работа)}

Подстановка линейных определяющих соотношений в (20)--(21) даёт члены
самозацепления, квадратичные по полям $T^a,I^a$ с теми же модулями
$a_1,f_1$, без новой свободной константы:
\begin{equation}
\pi_a=(e_a\lrcorner T^b)\wedge f_1(\star I_b),\qquad
E_a = e_a\lrcorner\Big[\tfrac12 f_1 I\wedge\star I-\tfrac12 a_1 T\wedge\star T\Big]
+a_1(e_a\lrcorner T^b)\wedge\star T_b - f_1(e_a\lrcorner I^b)\star I_b.
\end{equation}
Для двух источников ($T=T_1+T_2$, $I=I_1+I_2$) перекрёстные члены
$T_1\wedge\star I_2+T_2\wedge\star I_1$ дают $R$-зависимое
взаимодействие. При переносе на точечные прецессирующие монополи
(аналогия Брагинского, не точное тождество) и использовании
$I^a\sim (v/c)\,T^a$ (поскольку $I^a$ строится из градиента скорости)
получена оценка:
\begin{equation}
\langle I_{TT}(R,\Delta\phi)\rangle = \frac{\pi p^2}{R}\cos\Delta\phi,
\qquad p=2q_m a,
\end{equation}
проверенная численно методом Монте-Карло с точностью около 5\%, с
дополнительным подавляющим множителем $(v/c)\cdot f_1$ вместо
свободной константы $\lambda$ прежней модельной оценки.

\subsection{Сводная таблица}

\begin{center}
\begin{tabular}{|l|c|c|c|}
\hline
Вариант & Ведущий порядок & Знак/$\Delta\phi$? & Кулон $R^{-1}$ со знаком? \\
\hline
1. Линейная магнитостатика & $R^{-3}$ & да & Нет \\
2. Линейная плёнка Ферми & $R^{-1}$ & нет & Нет (без знака) \\
3. Нелинейная самосвязь & $R^{-1}$ & да & Да, форма верна; $f_1$ не откалиброван \\
\hline
\end{tabular}
\end{center}

\section{Заключение}

Явный вывод из точных полевых уравнений Лазара показывает, что
самозацепление калибровочного поля дислокаций структурно способно
породить знаковое, $R^{-1}$-взаимодействие между двумя электрически
нейтральными (по суммарному магнитному заряду) прецессирующими
диполями -- то есть верную функциональную форму электрического
кулона. В отличие от предыдущей модельной оценки, это не требует
введения новой, произвольно подбираемой константы связи: величина
эффекта фиксируется теми же упругими модулями $a_1,f_1$, что уже
определяют линейную (Пич--Кёлер) часть теории. Однако численное
значение этих модулей для физического вакуума -- в интерпретации
Дроздова, для самой упругой плёнки границы действия -- остаётся
неизвестным и не выводится из рассмотренных материалов; свобода
произвола, ранее заключённая в подборе показателя степени $R$, затем
в показателе степени массы, затем в величине $\lambda$, теперь
сосредоточена в физическом масштабе модулей $a_1,f_1$ вакуумной
плёнки. Это открытый вопрос для дальнейшей работы.

\begin{thebibliography}{9}
\bibitem{Fermi1923} E. Fermi, {\it О парадоксе массы в электродинамике},
Phys. Z. {\bf 24}, 340 (1923).
\bibitem{Lazar2010} M. Lazar, {\it On the Fundamentals of the
Three-dimensional Translation Gauge Theory of Dislocations},
Mathematics and Mechanics of Solids {\bf 16}(3), 253--264 (2010).
\bibitem{Braginsky2018} А.~Я.~Брагинский, {\it Вычисление сил,
действующих на импульс частицы в тензорном поле взаимодействия из
минимума действия. Объяснение аналогии между континуальной теорией
дислокаций и электродинамикой}, материалы конференции (2018).
\bibitem{Braginsky2016} А.~Я.~Брагинский, рабочие материалы (2016).
\bibitem{KadicEdelen1983} A. Kadi\'c, D.G.B. Edelen, {\it A Gauge
Theory of Dislocations and Disclinations}, Lecture Notes in Physics,
Vol. 174, Springer, Berlin (1983).
\end{thebibliography}

\end{document}
